% =============================================================================
%  CONCLUSION GÉNÉRALE ET PERSPECTIVES
% =============================================================================
\cleardoublepage
\entete{Conclusion générale}
\chapter*{CONCLUSION GÉNÉRALE}
\addcontentsline{toc}{chapter}{CONCLUSION GÉNÉRALE}

Au terme de ce travail, mené du 20 juin au 20 août 2026 au sein du \emph{Think
Tank} Foudjem de la Fondation Jean Félicien Gacha, il convient de revenir sur le
chemin parcouru et sur ce qu'il aura permis d'établir.

Le point de départ était un constat simple, formulé par l'institution
elle-même : la découverte de la Fondation supposait de s'y rendre. Cette
condition écartait les publics éloignés, la diaspora et les personnes à mobilité
réduite, et confinait un patrimoine considérable à son emprise géographique. Là
où une présence numérique existait, elle se réduisait à une vitrine figée :
photographies plates, notices brèves, objets isolés les uns des autres, aucune
possibilité d'interroger. La question posée était donc celle de la conception
d'une plateforme numérique unifiée capable d'offrir une visite en trois
dimensions, de relier l'institution à un réseau d'établissements pour retrouver
les œuvres apparentées, et d'assurer une assistance fiable fondée sur les seules
données de la Fondation.

Pour y répondre, nous avons conduit une recherche interactive de niveau
intégrateur, selon une approche mixte associant l'analyse documentaire,
l'observation directe, six entretiens semi-directifs et une enquête par
questionnaire auprès de quarante-deux visiteurs. La construction du système a
suivi le processus unifié 2TUP, modélisé en UML, et s'est appuyée sur une
démarche d'industrialisation continue.

Le travail a abouti à une plateforme opérationnelle, MUSÉA, articulant vingt
fonctionnalités autour de quatre ensembles : un progiciel de gestion couplé à un
site public éditable par l'institution ; une visite immersive à 360~degrés
complétée par la présentation en trois dimensions et la réalité augmentée à
taille réelle ; un assistant conversationnel ancré sur les documents validés par
la Fondation ; et un dispositif de rapprochement documentaire avec les
collections ouvertes du monde. L'ensemble repose sur une architecture cloisonnée
par la base de données, déployée sur une infrastructure en nuage dont le coût de
fonctionnement a été mesuré et maîtrisé.

Les quatre hypothèses spécifiques ont été vérifiées et l'hypothèse générale
corroborée, sous la réserve — explicitement énoncée — que l'institution produise
le contenu numérique dont le dispositif a besoin pour se déployer pleinement.

Au-delà de l'outil livré, ce travail dégage trois enseignements qui en
constituent la contribution propre.

\begin{enumerate}
  \item \textbf{Le décalage de vocabulaire constitue une seconde dispersion des
        œuvres.} Nous avons établi qu'une recherche menée en français dans les
        catalogues internationaux ramène du bruit là où la même recherche,
        reformulée dans le vocabulaire de catalogage de ces institutions,
        identifie des correspondances précises et vérifiables — le meilleur
        score passant de 57 à 85 sur 100. Une institution africaine qui
        interroge ces catalogues dans ses propres termes pourrait conclure à
        tort que ses œuvres sont sans équivalent connu. La dispersion physique
        des objets se double ainsi d'une dispersion documentaire, qui est,
        elle, réversible.

  \item \textbf{Dans une plateforme accueillant plusieurs organisations,
        l'isolement des données ne peut pas être une propriété du code.} Les
        deux fuites détectées durant le projet étaient l'une et l'autre
        d'origine applicative : le programme demandait sincèrement ce qu'il
        croyait devoir demander. Seule une garantie portée par la base de
        données elle-même résiste à l'erreur de programmation — et, tout aussi
        important, au détournement d'un agent automatique par une consigne
        malveillante.

  \item \textbf{La qualité d'une visite restituée repose sur quatre continuités
        — spatiale, d'échelle, documentaire et narrative — dont la rupture d'une
        seule suffit à faire échouer l'ensemble.} Ce modèle, dégagé des
        résultats, offre une grille de conception transposable à d'autres
        institutions.
\end{enumerate}

Sur le plan personnel, ce stage aura été l'occasion de conduire un projet
complet, de la formulation du besoin jusqu'à la mise en production, dans un
cadre professionnel réel et sous des contraintes qui ont obligé à écrire
soi-même ce que l'on aurait ordinairement importé. Cette contrainte, d'abord
subie, s'est révélée formatrice : elle a imposé de comprendre les mécanismes au
lieu de les employer.

Il reste que ce travail ne prétend ni à l'exhaustivité ni à la généralisation
statistique. Conduit auprès d'une institution unique, sur un échantillon de
convenance et dans un temps contraint, il décrit un dispositif, en mesure le
comportement et en tire des enseignements dont la transposition demandera
vérification. C'est précisément l'objet des perspectives qui suivent.

% =============================================================================
\cleardoublepage
\entete{Perspectives}
\chapter*{PERSPECTIVES}
\addcontentsline{toc}{chapter}{PERSPECTIVES}

Le dispositif construit ouvre un ensemble de prolongements, que nous présentons
regroupés en cinq axes, du plus immédiat au plus prospectif. Certains d'entre
eux n'ont pu être menés à terme durant le stage en raison de contraintes de
temps, de moyens ou d'accès ; nous le signalons explicitement le cas échéant.

\section*{Axe 1 — Consolider ce qui existe}
\addcontentsline{toc}{section}{Axe 1 — Consolider ce qui existe}

\begin{enumerate}
  \item \textbf{Comparer les œuvres sur l'image et non sur le texte.} C'est la
        perspective la plus importante. Les rapprochements documentaires
        reposent actuellement sur la description écrite des œuvres, faute de
        processeur graphique et de possibilité d'installer un modèle de vision
        sur le poste de travail. Le recours à un modèle de plongement visuel
        permettrait de rapprocher deux objets visuellement proches mais décrits
        selon des conventions différentes — situation fréquente entre un
        catalogue africain et un catalogue européen. Cette évolution supprimerait
        la principale limite de la mémoire réunifiée.

  \item \textbf{Produire les cartes de profondeur en production.} Le moteur de
        relief a été vérifié sur une carte fabriquée pour l'essai, mais aucune
        carte n'a encore été produite sur des œuvres réelles : le traitement
        exige des bibliothèques qui n'ont pu être installées. Une exécution dans
        un conteneur, telle que prévue par l'intervention proposée, lèverait cet
        obstacle.

  \item \textbf{Élargir l'angle de parallaxe du relief.} L'illusion est
        actuellement bornée à vingt degrés, au-delà desquels les zones absentes
        de la photographie s'étirent. Une passe de reconstitution des zones
        masquées permettrait d'élargir cet angle.

  \item \textbf{Obtenir les clés d'accès aux quatre collections écartées.}
        Europeana, le Rijksmuseum, la Smithsonian Institution et les Harvard Art
        Museums représentent un volume documentaire considérable. Leur
        intégration est déjà prévue techniquement ; elle ne dépend que d'une
        démarche administrative.

  \item \textbf{Éprouver la chaîne de rapprochement de bout en bout.} Chaque
        maillon a été vérifié isolément ; l'enchaînement complet reste à
        exécuter avec un compte du personnel.

  \item \textbf{Étendre les essais de réalité augmentée.} Le décodage des codes
        QR par caméra et le démarrage effectif d'une session n'ont pu être
        éprouvés que sur un nombre restreint d'appareils. Une campagne d'essais
        sur un parc représentatif s'impose.
\end{enumerate}

\section*{Axe 2 — Enrichir l'expérience du visiteur}
\addcontentsline{toc}{section}{Axe 2 — Enrichir l'expérience du visiteur}

\begin{enumerate}
  \item \textbf{Le mur de la dispersion.} L'écran de fin de parcours existe ;
        il reste à y animer le globe pour annoncer au visiteur : « les œuvres
        apparentées à celles que vous venez de voir se trouvent dans tel nombre
        de pays ». Le geste narratif serait fort : la visite s'achèverait sur la
        mesure de ce dont l'institution a été séparée.

  \item \textbf{Un guide vocal contextuel.} L'assistant sait déjà où se trouve
        le visiteur. Il pourrait commenter spontanément la salle où celui-ci
        vient d'entrer, à la manière d'un guide qui prend la parole sans qu'on
        l'interroge.

  \item \textbf{Le mode hors connexion.} L'application est déjà installable ; il
        resterait à mettre en cache les panoramas et les modèles des salles
        visitées, afin qu'une visite entamée puisse se poursuivre sur une
        connexion défaillante — situation courante pour une partie du public
        visé.

  \item \textbf{La reconnaissance d'une œuvre par photographie.} Un visiteur
        présent sur place pourrait photographier une pièce et obtenir
        immédiatement sa notice, son récit et ses œuvres apparentées, sans avoir
        à la chercher dans un catalogue.

  \item \textbf{La médiation ludique à destination des scolaires.} Les
        mécanismes de parcours à étapes existent déjà dans la base ; ils
        pourraient soutenir des jeux de piste conçus pour les classes, en
        présence comme à distance.

  \item \textbf{Les casques de réalité virtuelle.} Écartés à dessein pour des
        raisons d'accessibilité, ils pourraient être proposés en complément
        pour un usage sur place ou en médiathèque, l'interface du navigateur les
        prenant nativement en charge.
\end{enumerate}

\section*{Axe 3 — Ancrer le dispositif dans la culture locale}
\addcontentsline{toc}{section}{Axe 3 — Ancrer le dispositif dans la culture locale}

\begin{enumerate}
  \item \textbf{Le multilinguisme étendu aux langues du territoire.}
        L'application est aujourd'hui bilingue français-anglais. L'ajout des
        langues de l'Ouest camerounais — ghomálá', fe'efe'e, medumba —
        transformerait la portée du projet : il ne s'agirait plus seulement de
        faire connaître un patrimoine à l'extérieur, mais de le restituer aux
        communautés qui en sont issues, dans leur langue. La structure
        d'internationalisation du système le permet sans modification.

  \item \textbf{La collecte des récits oraux.} Le patrimoine immatériel — les
        récits, les usages, les protocoles associés aux objets — reste peu
        documenté. Une campagne d'enregistrement auprès des détenteurs de la
        mémoire, rattachée aux fiches d'œuvres et aux personnages de la
        généalogie, enrichirait considérablement la base documentaire de
        l'assistant.

  \item \textbf{L'approfondissement de la généalogie.} L'arbre fonctionne et se
        parcourt dans les deux sens ; il gagnerait à intégrer les migrations
        historiques, les alliances entre chefferies et les corpus d'objets
        rattachés à chaque règne.
\end{enumerate}

\section*{Axe 4 — Industrialiser et pérenniser}
\addcontentsline{toc}{section}{Axe 4 — Industrialiser et pérenniser}

\begin{enumerate}
  \item \textbf{Le passage aux conteneurs orchestrés}, décrit au chapitre~4,
        constitue le préalable à l'automatisation de l'ingestion documentaire,
        de la reconstruction photogrammétrique et de l'interrogation programmée
        des collections.

  \item \textbf{La migration vers Kubernetes} le jour où la plateforme dépassera
        une dizaine d'institutions ou exigera plusieurs services fonctionnant en
        permanence. La conteneurisation la rend possible sans réécriture.

  \item \textbf{Le renforcement du suivi de la dépense.} L'étiquetage des
        ressources par institution permettrait d'attribuer précisément le coût
        et, à terme, de fonder une tarification équitable de la plateforme.

  \item \textbf{La sauvegarde et la reprise après sinistre.} Les données
        patrimoniales sont irremplaçables : une politique de sauvegarde
        externalisée et un exercice de restauration périodique devraient être
        formalisés.

  \item \textbf{L'évaluation d'impact.} Une mesure de l'audience réelle — origine
        géographique des visiteurs, durée des parcours, taux d'adhésion,
        retombées sur la fréquentation physique — permettrait d'objectiver
        l'apport du dispositif au-delà de sa seule performance technique.
\end{enumerate}

\section*{Axe 5 — Ouvrir et essaimer}
\addcontentsline{toc}{section}{Axe 5 — Ouvrir et essaimer}

\begin{enumerate}
  \item \textbf{L'accueil d'institutions partenaires.} L'ouverture de
        l'architecture ne sera véritablement établie qu'après avoir été éprouvée
        par deux ou trois institutions distinctes de la Fondation.

  \item \textbf{L'ouverture des données de la Fondation.} La plateforme
        interroge aujourd'hui les catalogues du monde ; elle pourrait à son tour
        exposer les siens selon les mêmes standards, afin que les œuvres de
        Bangoulap deviennent trouvables depuis l'extérieur. Le mouvement serait
        alors réciproque, et non plus à sens unique.

  \item \textbf{Le dialogue avec les institutions détentrices.} Les
        rapprochements validés constituent une matière documentaire nouvelle,
        susceptible d'alimenter les échanges scientifiques avec les musées
        conservant des pièces apparentées — voire les discussions relatives à
        leur circulation.

  \item \textbf{Le modèle économique de la plateforme.} La mutualisation du
        catalogue mondial fait décroître le coût par institution à mesure que le
        réseau s'étoffe. Cette propriété mérite d'être instruite comme la base
        d'un modèle soutenable, permettant à de petites structures d'accéder à
        un outil qu'elles ne pourraient financer isolément.
\end{enumerate}
